\section{Defining Security}
\label{sec:threatmodel}

\noindent\textbf{Previous Work}
Previous work looking at enclaves in a context of transient execution attacks \cite{bourgeat2019mi6}
enforces \emph{Strong Microarchitectural Isolation} (SMI) by completely isolating security domains at the microarchitectural level.
See Section \ref{sec:hardwareisolation} on hardware isolation.
Unfortunately, this isolation forbids any shared architectural state between an enclave and the OS, disallowing critical functionalities such as shared memory.
Any interaction between the enclave and the OS has to go through the SM and requires several (expensive) context switches, microarchitectural flushes and data copies from one domain to another.
In \cite{bourgeat2019mi6} enclave completion time is considered outside of the threat model. Generally,
completion time can and should be handled independently using traffic shaping (padding) techniques.

\noindent\textbf{Secure Shared Memory}
We relax SMI to allow some shared memory between the enclave and the OS.
Nevertheless, to prevent transient execution attacks, we require that the enclave software and the hardware together achieve Speculative Non-Interference (SNI) \cite{guarnieri2020spectector, guarnieri2021hardware, behnia2021speculative, fabian2022automatic, yu2019speculative, yu2019speculative2}, i.e., that a victim program running in an enclave should not leak more secrets when running speculatively using our mechanisms than if run with non-speculatively (sequential execution).
However, similar to previous work, we do not consider the timing of microarchitectural events to be visible to the attacker.
In summary, no timestamps are associated with the observed shared memory accesses and the associated microarchitectural events; only the traces of the corresponding addresses are exposed.
We consider that leakage through the timing of shared memory accesses is similar to leakage through completion time, and it should be addressed using orthogonal traffic-shaping techniques such as \cite{deutsch2022dagguise, wang2014timing, shafiee2015avoiding, zhou2017camouflage}. This is left for future work.


\subsection{Attacker Capabilities}

We now complete our threat model by describing the attacker's capabilities.
The attacker is collocated on the same machine as the victim enclave and can mount any software-based attack, including microarchitectural side channel and transient execution attacks.
The attacker has compromised the majority of the software stack, including other malicious enclaves and the operating system (OS), except for a limited trusted code base (TCB) including the secure bootloader, the SM (together 15K lines of code (LOC)) and the victim program running inside of the enclave.
The attacker also has control over the contents of shared memory.
For the rest of this paper, we refer to any code outside of the TCB as the OS.
We ignore attacks requiring physical access to the machine such as physical (e.g., noise, electromagnetic field, etc.) side channels. 
We also ignore side channels that can be mounted in software, but exploit the physical layer like power side channels.
We do not consider Rowhammer attacks.
As highlighted by our security definition, we are worried about enclave privacy and do not consider denial of service.
We also do not protect enclave programs that explicitly leak their own secrets, including through shared memory.